\documentclass[11pt]{article}

\usepackage[a4paper,margin=1in]{geometry}
\usepackage{amsmath,amssymb,amsfonts}
\usepackage{booktabs}
\usepackage{longtable}
\usepackage{hyperref}
\usepackage{graphicx}
\usepackage{physics}

\title{MATC v16.9\\
A Reduced Einstein--Cartan--Proca Effective Theory\\
for Nonsingular Bianchi-IX Cosmology}
\author{}
\date{}

\begin{document}

\maketitle

\begin{abstract}
We investigate a reduced Einstein--Cartan--Proca (ECP) effective field theory describing a propagating axial-torsion sector coupled to a contracting Bianchi-IX cosmology. The framework is formulated as a metric-affine effective theory in which the irreducible axial component of torsion behaves as a massive pseudovector field with parity-violating derivative interactions. The primary objective is not the construction of an exact isotropic bounce, but the suppression of chaotic Belinskii--Khalatnikov--Lifshitz (BKL) oscillations through nonlinear axial backreaction.
\end{abstract}

\section{Introduction}

The problem of anisotropic instability remains one of the central obstacles facing nonsingular cosmology. In contracting spacetimes, anisotropic shear generically scales as

\[
\sigma^2 \propto a^{-6}
\]

leading to Mixmaster instability characteristic of BKL evolution.

This work investigates whether a propagating axial component of spacetime torsion may dynamically absorb anisotropic energy during contraction.

\section{Effective Einstein--Cartan--Proca Construction}

\subsection{Metric-Affine Origin}

The torsion tensor is defined by

\[
T^\lambda_{\mu\nu}
=
\Gamma^\lambda_{\mu\nu}
-
\Gamma^\lambda_{\nu\mu}
\]

The propagating degree of freedom retained in the EFT is the axial pseudovector

\[
A^\mu
=
\frac{1}{6}
\epsilon^{\mu\nu\rho\sigma}
T_{\nu\rho\sigma}
\]

\subsection{Effective Action}

The reduced action is

\[
S
=
\int d^4x \sqrt{-g}
\left[
\frac{M_{Pl}^2}{2}R
-\frac14 F_{\mu\nu}F^{\mu\nu}
+\frac12 M_A^2 A_\mu A^\mu
-\frac12 (\partial\phi)^2
-V(\phi)
+\frac{\alpha}{M_*}\phi F_{\mu\nu}\tilde F^{\mu\nu}
\right]
\]

with

\[
F_{\mu\nu}
=
\nabla_\mu A_\nu
-
\nabla_\nu A_\mu
\]

and

\[
\tilde F^{\mu\nu}
=
\frac12
\epsilon^{\mu\nu\rho\sigma}
F_{\rho\sigma}
\]

\section{Constraint Structure and Degrees of Freedom}

The total propagating count is:

\begin{itemize}
\item 2 tensor gravitational modes,
\item 3 Proca modes,
\item 1 scalar mode.
\end{itemize}

The canonical momentum conjugate to \(A_0\) vanishes:

\[
\pi^0 \approx 0
\]

Consistency generates the secondary constraint

\[
\partial_i \pi^i + M_A^2 A_0 \approx 0
\]

which fixes \(A_0\) algebraically.

\section{Covariant Energy Transfer Mechanism}

Variation with respect to \(A_\mu\) yields

\[
\nabla_\nu F^{\nu\mu}
+
M_A^2 A^\mu
=
-
\frac{2\alpha}{M_*}
\tilde F^{\mu\nu}
\partial_\nu \phi
\]

In a homogeneous cosmological background,

\[
\partial_\nu \phi = (\dot\phi,0,0,0)
\]

so only \(\tilde F^{\mu 0}\) contributes dynamically.

\subsection{Bianchi-IX Curvature Source}

The vector field is expanded in the invariant Bianchi-IX basis

\[
A = A_i(t)\omega^i
\]

where the Maurer--Cartan one-forms satisfy

\[
d\omega^i
=
-\frac12
C^i_{jk}
\omega^j \wedge \omega^k
\]

Consequently,

\[
F\tilde F \neq 0
\]

even in a homogeneous truncation.

\section{Nonlinear Damping and Scaling Behavior}

The interaction term modifies the continuity equation and may produce enhanced scaling:

\[
\rho_A \propto a^{-(6+\delta)}
\]

with

\[
\delta > 0
\]

emerging dynamically from scalar-to-axial energy transfer.

The viable region in parameter space is bounded by

\[
\alpha_c < \alpha < \alpha_{\max}
\]

where under-coupling fails to suppress BKL growth and over-coupling destabilizes the PDE structure.

\section{Hyperbolicity and Principal Symbol Analysis}

The most dangerous instability occurs in the longitudinal Proca mode. In the high-frequency WKB limit,

\[
v_L^2
\sim
1
-
\frac{\alpha^2 \dot\phi^2}
{M_*^2 M_A^2}
\]

Strong hyperbolicity requires

\[
v_L^2 > 0
\]

throughout the evolution.

\section{Numerical Program}

The reduced minisuperspace system evolves

\[
\mathbf U
=
(\Omega,\beta_+,\beta_-,A_i,\dot A_i,\phi,\dot\phi)
\]

using a fourth-order Runge--Kutta integrator.

\subsection{Constraint-Preserving Initialization}

Initial data satisfy

\[
\Omega_A(t_0) \ll 1
\]

with representative scales

\begin{center}
\begin{tabular}{ll}
\toprule
Quantity & Typical Scale \\
\midrule
$A_i(t_0)$ & $10^{-8} - 10^{-6} M_{Pl}$ \\
$\dot A_i(t_0)$ & $10^{-8} - 10^{-6} M_{Pl}^2$ \\
$\Omega_{shear}(t_0)$ & $0.8 - 0.95$ \\
\bottomrule
\end{tabular}
\end{center}

\subsection{Adiabatic Coupling Activation}

To reduce numerical stiffness,

\[
\alpha_{\mathrm{eff}}(t)
=
\alpha
\left(
1-e^{-(|H|/H_c)^n}
\right)
\]

with \(n \sim 2\).

\section{Numerical Diagnostics}

The numerical program succeeds only if:

\begin{enumerate}
\item BKL oscillations are suppressed,
\item enhanced scaling satisfies
\[
-\frac{d\ln\rho_A}{d\ln a} > 6
\]
\item the Hamiltonian residual
\[
\mathcal R_H
=
\left|
\frac{3H^2-\rho_{tot}}{\rho_{tot}}
\right|
\]
remains below \(10^{-6}\),
\item the Kretschmann scalar remains finite.
\end{enumerate}

\section{Failure Modes and Falsifiability}

\begin{center}
\begin{tabular}{lll}
\toprule
Failure Regime & Diagnostic & Interpretation \\
\midrule
Under-coupled & $\delta \le 0$ & BKL collapse resumes \\
Over-coupled & $v_L^2 < 0$ & Hyperbolicity breakdown \\
EFT Breakdown & Higher operators dominate & EFT invalid \\
Constraint Failure & $\mathcal R_H \gg 10^{-6}$ & Numerical inconsistency \\
\bottomrule
\end{tabular}
\end{center}

\section{Conclusion}

We have presented a reduced Einstein--Cartan--Proca effective theory describing a propagating axial-torsion sector coupled to a contracting Bianchi-IX cosmology.

The proposal stands or falls entirely on computational verification.

If the reduced evolution system demonstrates:

\begin{itemize}
\item enhanced axial scaling,
\item bounded anisotropy,
\item preserved Hamiltonian constraints,
\item finite curvature,
\end{itemize}

then the framework constitutes a technically viable nonsingular modified-gravity EFT.

\end{document}